\section{Intrinsic Smith Data and Presentations}
\label{sec:intrinsic-smith}

An executable SNF result contains transformation matrices, inverse
certificates, and a canonical matrix. Applications usually need the diagonal
invariants without the execution trace. The public structure
\code{SmithData} records that smaller interface:

\begin{lstlisting}[style=leanstyle,caption={Ordered intrinsic Smith data.}]
structure SmithData (m n : Nat) (R) where
  rank           : Nat
  rank_le        : rank <= min m n
  diagonal       : Fin (min m n) -> R
  nonzero_prefix : forall i, i.1 < rank -> diagonal i != 0
  zero_tail      : forall i, rank <= i.1 -> diagonal i = 0
  normalized     : forall i, diagonal i = normalize (diagonal i)
  divisibility   : forall i j, i.1 <= j.1 ->
    Dvd.dvd (diagonal i) (diagonal j)
\end{lstlisting}

The diagonal has length \(\min(m,n)\). It retains unit entries and every zero
after the rank. \code{SNFResultFin.smithData} reads the structure from the
canonical matrix, while the arbitrary-index result reads the same stored
\code{Fin} matrix. Changing the explicit indexing therefore leaves the data
unchanged.

\subsection{Signatures and determinantal ideals}

\code{SmithSignature} forgets the order and chosen representatives. It stores
the row and column counts, rank, and a multiset of \code{Associates R}
containing all nonzero factors. Unit factors remain in the multiset; the zero
tail does not. The signature derives right-kernel rank
\(\mathit{cols}-\mathit{rank}\), left-kernel rank
\(\mathit{rows}-\mathit{rank}\), and cokernel free rank
\(\mathit{rows}-\mathit{rank}\).

For a \code{Fin}-indexed matrix \(A\), let \(I_k(A)\) be the ideal generated by
its \(k\times k\) minors. The development proves
\[
 I_0(A)=\top,\qquad
 I_k(S)=\left(\prod_{i<k} d_i\right)
   \quad (1\leq k\leq r),\qquad
 I_k(S)=\bot\quad(k>r)
\]
when \(S\) is a Smith matrix with rank \(r\) and diagonal \(d_i\), matching the
classical relation between determinantal divisors and invariant factors
\citep[p. 81]{newman1971smith}. The
oversized-minor theorem also gives \(I_k(A)=\bot\) for
\(k>\min(m,n)\).

Left and right multiplication give ideal inclusions. Applying the inverse
certificates gives the reverse inclusions, so
\code{determinantalIdealFin_eq_of_certificate} proves equality for every
strong two-sided certificate. Explicit permutation matrices prove the same
invariance for row and column reindexing.

Equality of all determinantal ideals determines the first index at which the
ideals vanish and hence determines the rank. For indices inside the nonzero
prefix, equality of consecutive prefix-product ideals recovers each factor up
to associates. The resulting theorems compare rectangular Smith data with
different zero-tail dimensions and conclude equality of their nonzero
associate-factor multisets. This yields an intrinsic signature independent of
the algorithm, transformation matrices, and finite indexing.

\subsection{The mathlib PID bridge}

At the manifest-pinned mathlib commit, the PID API
\citep{mathlib4320pid} exposes
\code{Module.Basis.SmithNormalForm}, a valid basis
decomposition, but its coefficient witness does not include the canonical
normalization and divisibility chain needed for exact executable-factor
comparison. The project wraps such a witness in
\code{CanonicalPIDSmithBasis}, whose extra field proves that the associated
diagonal matrix satisfies \code{IsSmithNormalFormFin}.

The matrix input and this diagonal basis have equal determinantal ideals.
Intrinsic uniqueness then gives three public comparison levels:

\begin{itemize}
  \item equal multisets of associate classes over the Euclidean/PID layer;
  \item equal multisets after applying canonical normalization;
  \item equal sorted \code{Int.natAbs} lists for integer matrices.
\end{itemize}

No theorem equates the raw mathlib coefficients with executable factors at
matching indices. Such a statement would depend on choices that the raw
witness does not canonicalize.

\subsection{Column-oriented finite-free presentations}

The presentation interface fixes one matrix direction:

\begin{lstlisting}[style=leanstyle,caption={Finite-free presentation seam.}]
structure FiniteFreePresentation (R) where
  numGenerators : Nat
  numRelations  : Nat
  relationMatrix :
    Matrix (Fin numGenerators) (Fin numRelations) R
\end{lstlisting}

The matrix denotes
\[
  R^{\mathit{numRelations}}
    \longrightarrow R^{\mathit{numGenerators}}
\]
through \code{Matrix.mulVecLin}. Its columns are relations, and
\code{presentedModule} is the cokernel. The Smith result supplies a linear
equivalence
\[
  \operatorname{coker}(A)
  \simeq
  \left(\prod_{i<r} R/(d_i)\right)
    \times R^{m-r}.
\]
This theorem-facing product retains unit factors, whose quotient modules are
zero. A separate display reader filters unit factors.

The presentation layer defines
\(\operatorname{Fitt}_k=\smash{I_{m-k}(A)}\), using the standard convention in
which Fitting ideals are generated by minors of a presentation
\citep[Definition 15.8.3]{stacksFittingIdeals}. It records the usual top ideal
once \(k\geq m\) and the zero ideal when the requested minors exceed the
matrix dimensions. The executable Smith kernel does not enumerate minors;
Fitting ideals serve the semantic presentation theory.

Each matrix column also becomes one finitely supported
\code{Module.Relations.relation}. The project proves that mathlib's quotient
by their span is linearly equivalent to \code{presentedModule}, then packages
the equivalence as a \code{Module.Presentation}. The existing abelian-group
API now supplies the integer specialization. A standalone patch records the
generic matrix-relations bridge for possible upstream use; the pinned mathlib
dependency does not apply that patch.
